\section{Experimental Setup}
\label{sec:setup}
% EDIT (2026-10-29): new section built from the old 4.1--4.2 (workloads, models, hardware), 7.1
% (applications table) and the setup paragraphs of Sections 5--8; one subsection per experiment, in
% the order the experiments were designed. The measurement protocol is stated in Section 2 and only
% referenced here.
We ran five experiments on the loop of Section~\ref{sec:method}, designed in the order they are
numbered: a prompt-component ablation on single-kernel benchmarks (E1), a test of the
semantic-preservation constraint on the benchmark that was gamed (E2), the end-to-end study on two
production applications (E3), an ablation that removes the profiling data (E4), and a comparison of
raw against interpreted profiling data (E5). This section states the workloads, models and hardware
they share and then defines each experiment.

\subsection{Benchmarks and Applications}
\label{sec:explanation}
\paragraph{Benchmarks.}
The prompt-component ablation uses six HeCBench \cite{jin2023hecbench} CUDA benchmarks---\texttt{atan2},
\texttt{bitonic-sort}, \texttt{md5hash}, \texttt{accuracy}, \texttt{sobol} and
\texttt{shmembench}---chosen to span memory-bound (\texttt{shmembench}) and compute-bound
(\texttt{bitonic-sort}, \texttt{md5hash}, \texttt{atan2}) behaviour and to include kernels with
built-in verification. The later benchmark experiments use a further ten HeCBench programs:
\texttt{adv}, \texttt{amgmk}, \texttt{atomicSystemWide}, \texttt{background-subtract},
\texttt{bincount}, \texttt{blockAccess}, \texttt{boxfilter}, \texttt{ccsd-trpdrv}, \texttt{concat} and
\texttt{convolutionSeparable}. An eleventh, \texttt{chemv}, was run but is excluded from every table
because its built-in check cannot pass for the unmodified program (Section~\ref{sec:guidvar}).

\paragraph{Applications.}
The end-to-end study uses two production codes: AutoDock-GPU (molecular docking, CUDA) and Cholla
(astrophysical hydrodynamics, CUDA). Each is built from its own build system with its own inputs and
validated by the gate of Section~\ref{sec:correctness}. For each, we profiled the whole application
with Nsight Systems and Nsight Compute, selected the source file containing the kernel that dominates
GPU time---after confirming that this kernel also dominates \emph{application} time---and gave the
pipeline that single file to optimize, exactly as for a benchmark (Table~\ref{tab:apps}).

\begin{table}[t]
\centering
\footnotesize
\caption{Production applications, the file given to the pipeline, and the share of GPU time held by
  the target kernel.}
\label{tab:apps}
\setlength{\tabcolsep}{3pt}
\begin{tabular}{@{}llr@{}}
\toprule
\textbf{Application} & \textbf{Target file (hot kernel)} & \textbf{Share} \\
\midrule
AutoDock-GPU & \texttt{calcMergeEneGra.cu} & 98.7\% \\
Cholla & \texttt{hllc\_cuda.cu} & 56.0\% \\
\bottomrule
\end{tabular}
\end{table}

\subsection{Models}
\label{sec:models}
% EDIT (2026-10-29): the reason for the model change is the author's statement (review of
% 2026-10-29); no document in the experiments repository records the date it was disabled.
The prompt-component ablation (E1) uses GPT-5 and Gemini~2.5~Pro. Gemini~2.5~Pro was disabled on the
API we use after that study, so the four later experiments use its successor, Gemini~3.1~Pro, together
with GPT-5. We note the change wherever numbers from E1 appear next to numbers from the later
experiments, since the two Gemini versions are not directly comparable.
% Earlier exploratory runs used llama3.1:70B, llama2:7B and qwen-coder:14B. They are not in any
% reported table; add a paragraph here if you want to include those numbers.

\subsection{Hardware and Software}
\label{sec:hardware}
E1, E2 and the benchmark half of E4 are timed on an NVIDIA H100 PCIe (sm\_90); E3, the application
half of E4 and E5 are timed on an NVIDIA GH200 (sm\_90). Both nodes use CUDA 12.6.3, except the
benchmark half of E4, which was rebuilt and re-timed with CUDA 13.3.1 after a driver update on the H100
node made binaries built with CUDA 12.x produce wrong output there (Section~\ref{sec:limits}). The
profiles supplied in the prompt for the six E1 benchmarks were collected on an A100 ($108$ SMs) rather
than on the H100 ($114$ SMs) they were then timed on; the ten later benchmarks use native H100
profiles, and the Cholla profile was collected on the GH200 on which it was measured. Both nodes are
shared with other users, which is why every timing follows the protocol of Section~\ref{sec:timing}.

\subsection{Experiment 1: Prompt-Component Ablation}
\label{sec:setup-e1}
E1 answers RQ1. It runs the six variants of Section~\ref{sec:ablation} with both models on the six
benchmarks: $2 \times 6 \times 6 = 72$ runs, each with up to three repair iterations. Every run
supplies PC-sampling and roofline summaries, so the variants differ only in the instructions. We record
for each run whether it produced a valid candidate, at which iteration, and the candidate's runtime,
which for this experiment is the in-loop ten-run average (Section~\ref{sec:timing}); a run is
\emph{valid} when its candidate built and passed the benchmark's check. Baseline runtimes of the
unmodified program were recorded for \texttt{bitonic-sort}, \texttt{sobol} and \texttt{shmembench},
so speedups are stated for those three and the other three benchmarks compare variants with one
another only.

\subsection{Experiment 2: The Semantic-Preservation Constraint}
\label{sec:setup-e2}
E2 answers RQ2 on \texttt{atomicSystemWide}, the benchmark both models gamed. It compares the
unconstrained baseline prompt against three prompts that carry the semantic-preservation block of
Section~\ref{sec:constraint}: the relaxed \emph{structural} prompt, and the \texttt{pcrl} and
\texttt{code\_only} prompts, which keep the original prohibitions and add only the block. The
unconstrained and structural arms were generated and timed in one interleaved session on the H100 node;
the \texttt{pcrl} and \texttt{code\_only} arms were generated later and timed on the same node in the
re-timing of E4 (CUDA 13.3.1, five clean repetitions per arm). Both GPT-5 and Gemini~3.1~Pro are run.
Besides the timings, we read the generated source of every arm to see whether the per-iteration atomic
operation is still there.

\subsection{Experiment 3: Production Applications}
\label{sec:setup-e3}
E3 answers RQ3. On each application, both models are run with the \emph{baseline} prompt and with the
\emph{structural} prompt, and the results are compared with an expert hand-tuned reference restricted
to the same single file. Each (application, model, prompt) cell is one generation, and a cell is
reported as a failure when no candidate passed the gate within three repair iterations. Timing uses
the application's own phase timer on the GH200. The clean repetitions per median are: AutoDock-GPU
model arms $5$ (idle GPU), except Gemini \emph{structural}, which was timed in a session with another
user's job resident and whose baseline ratios reproduced the idle session's within $0.7\%$;
AutoDock-GPU hand reference $7$; Cholla $3$ of $5$ (two excluded for a foreign job, spread
$\le 0.12\%$).

\subsection{Experiment 4: Profile-Data Ablation}
\label{sec:setup-e4}
E4 answers the first half of RQ4: does profiling data help at all? It compares the \texttt{pcrl}
prompt against the \texttt{code\_only} prompt of Section~\ref{sec:ablation}, which differ only in the
presence of the profile summaries, with both models on the ten later benchmarks and on both
applications: $24$ (program, model) pairs. The benchmark pairs are timed on the H100 node with five
repetitions per arm, all five clean for every pair; per-repetition spread there ranges from $0.1\%$ to
$37.4\%$, so each pair is scored with the tie rule of Section~\ref{sec:timing}. The application pairs
are timed on the GH200: Cholla over $4$ of $7$ repetitions (three excluded for drift with no foreign
job detected), AutoDock-GPU over $2$ of $7$ (five excluded for foreign jobs). For the applications a
second, independently generated \texttt{pcrl} sample was also run, which is the baseline arm of E3.

\subsection{Experiment 5: Prescriptive Guidance}
\label{sec:setup-e5}
E5 answers the second half of RQ4: what happens when the profiling data is converted into named
remedies before the model sees it? It runs four arms on the GH200, with the node idle, all arms built
once and timed round-robin in a single session: the unmodified program, the constrained
\emph{baseline} prompt, the relaxed \emph{structural} prompt, and \emph{diagnostics}, which is the
structural prompt plus the checklist (Section~\ref{sec:ablation}). The cleanest comparison is
therefore \emph{diagnostics} against \emph{structural}, which differ only in the checklist. Seven
benchmarks were run with both models; with \texttt{chemv} excluded, twelve (benchmark, model) pairs
are scored: \texttt{amgmk}, \texttt{bincount}, \texttt{blockAccess}, \texttt{boxfilter},
\texttt{concat} and \texttt{convolutionSeparable}. Each cell is a single generation.
